\section{\texorpdfstring{Wild inertia at \(p=3\)}{Wild inertia at p=3}}
\label{sec:wild-three}

The prime $3$ is the only place where the filtered inertia is unique but the
decomposition group is not.  We retain every finite alternative until it is
removed by an explicit local-field condition.  This order is important:
Serre's conjugation law distinguishes the last two inertia embeddings, but
it does not produce the candidate list.

\subsection{Complete decomposition and inertia inventory}

The local degrees in the line field and the certified double-six degrees are
\[
 27:(3,6,9,9),\qquad 36:(3,3,3,9,18).
\]
We enumerate all $350$ conjugacy classes of subgroups in the table of marks
for \texttt{U4(2).2}, reconstruct their actions on both labelled carriers,
and retain precisely those with the two displayed orbit partitions.

\begin{proposition}[Raw \(p=3\) hits]
\label{prop:p3-hits}
The simultaneous orbit targets have exactly three subgroup hits:
\[
\begin{array}{c|c|c|c}
\text{ToM index}&|H|&\IdGroup(H)&\text{initial role}\\ \hline
140&18&(18,4)&D\text{ or }I\\
142&18&(18,3)&D\text{ or }I\\
206&36&(36,10)&D\text{ only}.
\end{array}
\]
The complete list compatible with inertia normality and a cyclic residue
quotient is
\begin{equation}
(D,I,|D/I|)=(140,140,1),(142,142,1),(206,140,2),(206,142,2).
\label{eq:p3-all-di}
\end{equation}
\end{proposition}

\begin{proof}
The simultaneous orbit scan is exhaustive over the $350$ classes, so it
gives the three raw hits without a subgroup-name guess.  Their $3$-Sylow
subgroup has order $9$.  If \Tom{206} were inertia, the quotient of order
$4$ by wild inertia would be $V_4$, whereas $I_0/I_1$ must be cyclic of
order prime to $3$.  Hence \Tom{206} cannot be $I_0$.  It contains normal
subgroups in each of the two order-$18$ classes, both with cyclic quotient
of order $2$.  Testing normality and the residue quotient for every ordered
containment gives exactly \cref{eq:p3-all-di}.
\end{proof}

The abstract group structures clarify the remaining ambiguity:
\[
 \Tom{140}\cong(C_3^2):C_2,\qquad
 \Tom{142}\cong C_3\times S_3.
\]
Both contain the required wild group $P=C_3^2$; they differ in how the tame
involution acts on the eventual deep subgroup.  The order-$36$ class has no
role as inertia anywhere below.

\subsection{\texorpdfstring{Exhausting the deep \(C_3\)}{Exhausting the deep C3}}

The four branches of $E\otimes\Q_3$ have target different vector
\begin{equation}
 \mathbf d=(3,7,18,18).
\label{eq:p3-target-different}
\end{equation}
The tame $I_0$ contribution is $\mathbf b=(2,5,8,8)$.  One lower layer equal
to $P=C_3^2$ contributes $\mathbf u=(1,2,4,4)$.  Let $r$ be the number of
$P$ layers and $s$ the number of layers equal to a chosen deep
$Q\cong C_3$.  For every candidate $Q$, the branch equations are
\begin{equation}
 \mathbf b+r\mathbf u+s\mathbf w_Q=\mathbf d,\qquad
 r,s\in\Z_{\geq0}.
\label{eq:p3-branch-equation}
\end{equation}

The deep group is not chosen by abstract isomorphism type: three different
$C_3$ embeddings occur in the valid pairs.  Their complete profile multiset
is \Tom{6} twice, \Tom{7} once, and \Tom{8} once.  The two occurrences of
\Tom{6} come from distinct ordered $(D,I)$ containments; neither is
suppressed before the different equation is solved.

\begin{table}[ht]
\centering
\caption{Exact deep-$C_3$ exhaustion.  Orbit notation is run-length
notation, and every fraction is exact.}
\label{tab:p3-deep-exhaustion}
\resizebox{\textwidth}{!}{%
\begin{tabular}{@{}c c c c c c c@{}}
\toprule
$Q$&multiplicity&line type&double-six type&
$\dim(V_6^Q,V_{20}^Q)$&$\mathbf w_Q$&formal $(r,s)$\\
\midrule
\Tom{6}&2&$3^9$&$3^{12}$&$(0,8)$&
$(\frac13,\frac23,1,1)$&$(7,-18)$\\
\Tom{7}&1&$1^9 3^6$&$1^6 3^{10}$&$(4,10)$&
$(0,0,1,1)$&$(1,6)$\\
\Tom{8}&1&$3^9$&$1^3 3^{11}$&$(2,6)$&
$(\frac13,\frac23,1,1)$&$(7,-18)$\\
\bottomrule
\end{tabular}}
\end{table}

\begin{proposition}[Unique branch-different solution]
\label{prop:p3-deep}
\Cref{eq:p3-branch-equation} has a nonnegative integral solution
only for $Q=\Tom{7}$, and that solution is $(r,s)=(1,6)$.  Consequently
\[
 I_1=C_3^2,\qquad I_2=\cdots=I_7=\Tom{7}\cong C_3,\qquad I_8=1.
\]
\end{proposition}

\begin{proof}
For \Tom{6} and \Tom{8}, the first and third components of
\cref{eq:p3-branch-equation} give
\[
 r+\frac{s}{3}=1,\qquad 4r+s=10,
\]
whose unique formal integral solution is $(7,-18)$.  It cannot be a
filtration length.  For \Tom{7}, the first component gives $r=1$, and the
third gives $8+4+s=18$, hence $s=6$.  The second and fourth components agree.
The candidate list in \cref{tab:p3-deep-exhaustion} is exhaustive, so no
fourth profile remains.
\end{proof}

The fractional entries in \cref{tab:p3-deep-exhaustion} are not rounded
orbit heuristics.  They are the exact normalized contributions
$|Q|/|I_0|$ to the four branch permutation conductors.  The negative solution
therefore rejects an embedding, not merely a proposed break number.

\subsection{The odd-grade tame action}

After \cref{prop:p3-deep}, the only remaining choice is how the tame quotient
$C_2=I_0/I_1$ acts on the selected $Q=G_7/G_8\cong C_3$.  Serre's graded
action law states that, for $s\in G_0$ and
$\tau\in G_i/G_{i+1}$,
\begin{equation}
 \theta_i(s\tau s^{-1})=\theta_0(s)^i\theta_i(\tau).
\label{eq:serre-graded-action}
\end{equation}
Here $\theta_0(s)=-1$ for the nontrivial tame involution.  The precise source
is \citet[Chapter IV, \S2, Proposition 9, pp.~69--70]
{serre1979local}.  It supplies \cref{eq:serre-graded-action}; it does not
identify any Table-of-Marks class for the present surface.

\begin{proposition}[Serre discriminator]
\label{prop:p3-serre}
The inertia subgroup is \Tom{140}.  The central embedding \Tom{142} is
impossible, and the final ordered pairs are exactly
\begin{equation}
 (D,I)=(140,140),\qquad(206,140).
\label{eq:p3-final-di}
\end{equation}
\end{proposition}

\begin{proof}
The last nonzero grade is $i=7$.  Substitution in
\cref{eq:serre-graded-action} gives
\[
 \theta_7(s\tau s^{-1})=(-1)^7\theta_7(\tau)=-\theta_7(\tau),
\]
so the tame involution acts by inversion on $Q\cong C_3$.  In \Tom{140},
the selected normal \Tom{7} is inverted.  In \Tom{142}, it is centralized.
The latter contradicts the graded action law.  Removing the two rows with
inertia $142$ from \cref{eq:p3-all-di} leaves precisely
\cref{eq:p3-final-di}.  The selected deep subgroup is normal in each
surviving decomposition group.
\end{proof}

Serre's proposition acts only after the $350$-class scan, the four valid
pairs, and the three exact deep-profile equations have been exhausted; it
does not by itself select class $140$.

\subsection{Character consequences and the unresolved overgroup}

The fixed dimensions for the three nontrivial layers are
\begin{equation}
\begin{array}{c|ccc}
H&I_0=\Tom{140}&P=C_3^2&Q=\Tom{7}\\ \hline
\dim V_6^H&0&0&4\\
\dim V_{20}^H&3&4&10.
\end{array}
\label{eq:p3-fixed-dimensions}
\end{equation}
Thus the tame codimensions are $(6,17)$, the $P$ codimensions are $(6,16)$,
and the $Q$ codimensions are $(2,10)$.  Direct substitution into
\cref{eq:swan-definition} yields
\begin{align}
 \Sw_3(V_6,V_{20})
 &=\frac9{18}(6,16)+6\frac3{18}(2,10)=(5,18),
\label{eq:p3-swan}\\
 a_3(V_6,V_{20})&=(6,17)+(5,18)=(11,35).
\label{eq:p3-artin}
\end{align}
Their sum $46$ equals the direct different exponent in
\cref{tab:local-rows}; \cref{sec:conductors} gives the general conductor
argument and branchwise checks.

The two possibilities in \cref{eq:p3-final-di} have orders $18$ and $36$.
Both contain the same filtered inertia and the same labelled fixed spaces,
so \cref{eq:p3-swan,eq:p3-artin} and both field discriminants are independent
of this choice.  A decomposition Frobenius or a bad Euler polynomial would
not be independent of decomposition data.  We make no such inference here,
and resolving $|D_3|$ alone would still not provide epsilon factors or root
numbers.
